\documentclass[11pt,a4paper]{article}

\usepackage{fontspec}
\usepackage[magyar]{babel}
\usepackage{amsmath,amssymb,amsthm}
\usepackage{hyperref}

\hypersetup{
  colorlinks=true,
  linkcolor=blue,
  urlcolor=blue,
  citecolor=blue
}

\newtheorem{tetel}{Tétel}[section]
\newtheorem{allitas}[tetel]{Állítás}
\newtheorem{lemma}[tetel]{Lemma}
\newtheorem{kovetkezmeny}[tetel]{Következmény}

\theoremstyle{definition}
\newtheorem{definicio}[tetel]{Definíció}
\newtheorem{pelda}[tetel]{Példa}
\newtheorem{megjegyzes}[tetel]{Megjegyzés}

\renewcommand{\proofname}{Bizonyítás}

\DeclareMathOperator{\Frac}{Frac}
\DeclareMathOperator{\Spec}{Spec}
\DeclareMathOperator{\ord}{ord}
\DeclareMathOperator{\length}{length}

\newcommand{\OO}{\mathcal O}
\newcommand{\mm}{\mathfrak m}
\newcommand{\pp}{\mathfrak p}
\newcommand{\qq}{\mathfrak q}
\newcommand{\ZZ}{\mathbb Z}
\newcommand{\NN}{\mathbb N}
\newcommand{\kk}{\Bbbk}

\title{Értékelések és értékelésgyűrűk\\
\large Különös tekintettel a diszkrét értékelésgyűrűkre}
\author{Előadásjegyzet az Algebrai Görbék c. tárgyhoz}
\date{}

\begin{document}





\section{Teljesség, Cohen struktúratétele}

\subsection{Mikor legyen teljes egy DVR?}

Egy testen definiálhatunk egy normát egy adott diszkrét értékelésből.

\begin{definicio}
Legyen $F$ test, $\nu: F^{\times}:\rightarrow\mathbb{Z}$ értékelés, és $c\in\mathbb{R}_{>1}$ rögzített szám. Ekkor minden $x\in F^{\times}$-nek definiáljuk a következő normáját:
$$||x||_{\nu}:=c^{\nu(x)}\in\mathbb{R}_{\ge0}.$$
Továbbá a $0$ normája természetes módon
$$||0||_{\nu}=c^{-\infty}=0.$$
\end{definicio}

Azért hívjuk normának, mert pontosan a $0$-n vesz fel $0$-t, és multiplikatív. Továbbá az alábbi állításban szereplő harmadik tulajdonság, az {\it ultrametrikus egyenlőtlenség} okán nemarkhimédeszi normának is hívjuk.

\begin{allitas}
A fenti $||\cdot ||_{\nu}$ egy {\it nemarkhimédeszi norma}, azaz minden $x,y\in F$-re
\begin{enumerate}
\item $||x||_{\nu}=0\Leftrightarrow x=0$,
\item $||xy||_{\nu}=||x||_{\nu}\cdot ||y||_{\nu}$,
\item $||x+y||_{nu}\le\max\{||x||_{\nu},||y||_{\nu}\}$.
\end{enumerate}
\end{allitas}
\begin{proof}
Az 1. állítás teljesen világos.

A 2. állításhoz felírjuk, hogy a diszkrét értékelés definíciója szerint $\nu(xy)=\nu(x)+\nu(y)$. Innen $c^{\nu(x)+\nu(y)}=c^{\nu(x)}c^{\nu(y)}$, és ez kellett.

A 3. állításhoz felírjuk, hogy a diszkrét értékelés definíciója szerint $\nu(x+y)\le\max\{\nu(x),\nu(y)\}$, majd használjuk a $c^{\cdot}$ függvény monotonitását (ami $c>1$ miatt igaz).
\end{proof}

Mivel az ultrametrikus egyenlőtlenségből következik a háromszög-egyenlőtlenség, ez a norma indukál egy metrikát, amivel $F$ metrikus tér lesz. Így definiálható $F$ teljessége, ebből pedig a megfelelő DVR teljessége is.

\begin{definicio}
Az $F$ {\it teljes} a $||\cdot ||_{\nu}$ normára nézve, ha a normára nézve minden Cauchy-sorozat konvergál. Az $\OO_{\nu}$ DVR {\it teljes}, ha $F$ teljes a $||\cdot ||_{\nu}$ normára nézve.
\end{definicio}

\subsection{Cohen struktúratétele}

Cohen struktúratétele azt írja le, hogy egy algebrai görbe lokálisan szinte mindenhol ugyanúgy néz ki.

\begin{tetel}[Cohen struktúratétele]
Legyen $A$ egy DVR, $m\triangleleft A$, és $k=A/m$ az $A$ maradékteste. Tegyük fel, hogy $k\subseteq A$, $k\hookrightarrow A\twoheadrightarrow A/m$ izomorfizmus, továbbá $A$ teljes. Ekkor $A\cong k[[t]]$.
\end{tetel}

Ezt most nem bizonyítjuk.










\end{document}
